% =====================================================================
% Section VI (Conclusion) of the 1015 draft, input from Papers/1015/main.tex after experiments.tex.
% Defines sec:conclusion. Limitations and future work appear only here.
% =====================================================================

\section{Conclusion}
\label{sec:conclusion}

This paper studied zero-shot scale generalization in multiagent target defense. We identified the threat context, the relation among attackers, their targets, and the teammates that cover those targets, as the information a defender must build before it acts. LEAF builds this context by looping one weight-shared entity-attention operator over the current entities. A query formed from the defender's current state reads every round, learned weights fuse the reads, and recurrent memory enters only after the read. Sampled-depth training makes every round count a usable policy.

Trained with at most 10 defenders and 3 targets and deployed without retraining to 50 defenders or 12 targets, LEAF had the lowest damage in every test configuration, and even its weakest training run outperformed the best run of every baseline. It also kept every target undamaged in more episodes and reduced the damage of the worst episodes. Single-factor ablations showed that context construction, multi-round refinement, and current-state reading each reduce damage, and the last two matter most under extrapolation. A baseline enlarged to the size of LEAF remained far behind, so the gain comes from how the computation is organized. A probe of the recurrent state showed that memory drifts with scale while the current self state does not, which explains why memory belongs after the read. Under sampled-depth training, LEAF placed about 99\% of the fusion weight on the first two rounds. A single round therefore kept nearly all of the protection at 29\% less decision time, whereas a model trained at a fixed depth collapsed when shortened. For automated defense systems, these results mean that one trained policy can serve teams of different sizes and onboard computers of different speeds.

Several limits of this study point to future work. All results come from simulation, in which defenders observe every entity without noise, the attackers follow a fixed rule, and each method was trained with three seeds. The tests reached 50 aircraft per side or 12 targets. Field deployment calls for onboard sensing with limited range, communication delays, and attackers that adapt their tactics, and it would benefit from evaluation against learned attackers. The execution depth is currently set by the operator, and a policy that chooses its number of rounds for each decision could save further computation. Extending the threat context to heterogeneous defenders and to targets of different value is another direction.
